% GDZ PPN236018647,native558 after83;HKUST559.
\NumberedItem{84}
\indent Classes dari in quovis ordine; hincque numerorum in terna quadrata discerpibilitas ad theoriam solidam reducta.
\par\noindent\hfill Brunsv[igae, 1798] Apr.

\medskip\begingroup\small
\indent \textsc{Gauss} hat hier wohl absichtlich \textls[120]{ordine} geschrieben statt, wie man zunächst erwarten würde, \textls[120]{genere}, weil im art. 287. der \textit{Disquisitiones arithmeticae}, Werke I, S. 336, der Nachweis gesondert für jede der beiden Ordnungen (proprie et improprie primitivi) geführt wird. Ganz genau müßte es heißen: Classes dari in quovis genere cuiusvis ordinis. Vergl. die Angabe Werke I, S. 476, Bemerkung zum art. 287. III der \textit{Disquisitiones arithmeticae}, »Demonstratione primum munita sunt Mense Aprili 1798«.
\par\noindent\hfill\textsc{Klein.}\quad\textsc{Bachmann.}
\par\endgroup

\NumberedItem{85}
\indent Demonstrationem genuinam compositionis virium eruimus.
\par\noindent\hfill Gotting[ae, 1798] Mai.

\medskip\begingroup\small
\indent Im Nachlaß haben sich keine auf diesen Beweis bezügliche Aufzeichnungen gefunden, aber in einem Briefe von F. L. \textsc{Wachter} an \textsc{Gauss}, datiert Altenburg, den 16. Dezember 1814, wird er erwähnt. Es heißt dort:

»\,.\ .\ . Den \textsc{Duchayla}schen Beweis für das Kräfteparallelogramm habe ich deswegen unrichtig genannt, weil ihm, in der Gestalt, die ihm \textsc{Poisson} in seiner Mechanik gegeben, die Beweiskraft mangelt. Ihr Beweis scheint mir noch wesentlich verschieden zu sein, weil er streng auf dem Begriff einer festen Ebene beruht, \textsc{Poisson} hingegen \textls[120]{ausdrücklich} nur auf den Satz sich stützt: daß es gleich sey, an welchem Punkte einer festen Linie eine Kraft wirkt, und dieser Satz scheint mir völlig unzureichend zu seyn .\ .\ .«.

Über \textsc{Wachter} vergl. P. \textsc{Stäckel}, Mathematische Annalen \textbf{54}, 1901, S. 49. \textsc{Duchaylas} wesentlich geometrischer Beweis findet sich in der Correspondance de l'école polytechnique, Nr. 4, Paris 1808; \textsc{Poisson} hat ihm in seinem \textit{Traité de mécanique}, Paris 1811, T. I, S. 11, eine analytische Gestalt gegeben.
\par\noindent\hfill\textsc{Stäckel.}
\par\endgroup
\par\bigskip\noindent X1.\hfill67\hspace*{.14\linewidth}
